\subsection{A larger finite list from radial polynomial weights}
\label{sec:radial-certificate}

\begin{proposition}[An exact two-moment interval certificate]
\label{prop:radial-certificate}
Among the $34$-element subsets $S\subseteq\{0,\ldots,63\}$, some class with
fixed values of $\sum_{a\in S}a$ and $\sum_{a\in S}\binom a2$ contains at least
\[
  5{,}133{,}798{,}314{,}667
\]
subsets. Consequently, for every prime $p>64$,
$C=\RS(\F_p,\{0,\ldots,63\},32)$ satisfies
$\ell_C(30/64)\ge5{,}133{,}798{,}314{,}667$, with each constructed
codeword agreeing in exactly $34$ coordinates.
\end{proposition}
\begin{proof}
Taking complements preserves class sizes, so work with $30$-subsets.
Put
\[
 X_a=2a-63,\qquad Y_a=\frac{(2a-63)^2-1365}{4},\qquad
 x=\frac12\sum_{a\in S}X_a,\quad y=\frac12\sum_{a\in S}Y_a.
\]
The pair $(x,y)\in\mathbb Z^2$ is an invertible affine transformation of the
two moment sums.  Set
\[
 T=341x^2+5y^2,\quad D=1{,}884{,}025,\quad
 Q=T/D,\quad A=\frac{851547}{580691}.
\]
The ancillary file \path{radial_certificates/certificate_degree12.json} specifies a rational
polynomial $h$ of degree $12$.  Use the degree-$25$ radial weight
\[
 w(Q)=(A-Q)h(Q/A)^2.
\]
It is nonpositive outside $Q<A$.  If $L$ denotes the largest class, then
\[
 \sum_{|S|=30}w(Q(S))
 \le L\sum_{(x,y)\in\mathbb Z^2}\max\{w(Q(x,y)),0\}.
\]
The numerator is evaluated exactly from the moments through degree $25$
of $T$, obtained by an integer dynamic program on the $32$ population pairs
$(u,v),(-u,v)$.  The denominator is an exact sum over the $210{,}135$ lattice
points with $T\le2{,}762{,}804$.  After clearing all coefficient denominators,
the ancillary file records the integer numerator and denominator, whose
ratio has ceiling $5{,}133{,}798{,}314{,}667$.
The standard-library script \path{radial_certificates/verify_publication.py} recomputes all
$676$ mixed moments, checks the saved polynomial and ratio, and independently
recomputes the denominator by summing expanded monomials over the full signed
lattice.  Numerical optimization was used only to discover the rational
polynomial; it is not an input to the verification.
Lemma~\ref{im:finite-list} converts the common-moment class into the
stated Reed--Solomon list. This improves the degree-$20$ conditional
certificate by approximately $1.17\%$; no optimality is claimed.
\end{proof}
